\section*{第 \ref{ch:b2:ore-deposits} 章——矿床}

\begin{solution}{exo:b2:ore-deposits:1}
(a) $\qty{0.5}{\percent} = \qty{5000}{ppm}$，且 $5000/28 = 179$。(b) $\qty{6}{g/t} = \qty{6000}{ng/g}$，
且 $6000/1.5 = 4000$。(c) 铬铁矿中铬的质量分数为 $2 \times 52.00/(55.85 + 104.00 + 64.00) = 0.4646$；
铬铁岩中为 $0.9 \times 0.4646 = 0.418$，即 \qty{418000}{ppm}，则
$418000/92 = 4540$。
\end{solution}

\begin{solution}{exo:b2:ore-deposits:2}
(a) $914 \times 0.0053 = \qty{4.84}{Mt}$ 铜，$914 \times 10^6 \times 0.40 = \num{3.66e8}$ 克，
即 \qty{366}{t} 金；$1640 \times 0.0044 = \qty{7.22}{Mt}$ 铜，\qty{738}{t} 金。(b) $1640 - 914 =
\qty{726}{Mt}$，其中含铜 $7.216 - 4.844 = \qty{2.37}{Mt}$，平均品位为
$2.372/726 = \qty{0.327}{\percent}$：多出的块体品位低于 \qty{0.3}{\percent} 的\omterm{def:b2:ore-deposits:grade}{边界品位}，
但又略高于它。(c) 它们回答的是两种\omterm{def:b2:ore-deposits:grade}{边界品位}下的问题；脱离\omterm{def:b2:ore-deposits:grade}{边界品位}，
二者都谈不上\omterm{def:b2:ore-deposits:deposit}{矿床}的规模。
\end{solution}

\begin{solution}{exo:b2:ore-deposits:3}
(a) $(19.3 - 1.00)/(2.65 - 1.00) = 18.3/1.65 = 11.1$。(b) 等速沉降要求 $(\rho_g - 1)d_g^2 =
(\rho_q - 1)d_q^2$，故 $d_g/d_q = \sqrt{1.65/18.3} = 0.300$：金颗粒的沉降速度与
约大三倍的石英颗粒相同。(c) $d_q = 0.1\sqrt{3.7/1.65} = \qty{0.150}{mm}$。
\end{solution}

\begin{solution}{exo:b2:ore-deposits:4}
以单位面积和单位密度计，盖层保留自身的铜，并接收淋滤柱中的铜：
$g = g_0(h_e + \eta h_l)/h_e$。(a) $0.4 \times (50 + 150)/50 = \qty{1.6}{\percent}$。
(b) $0.4 \times (50 + 0.8 \times 150)/50 = \qty{1.36}{\percent}$。(c) $h_l = 50 \times (1.7/0.4 - 1)
= \qty{163}{m}$。这些品位位于富集斑\omterm{def:b2:magma-emplacement:bodies}{岩盖}层的
0.4--\qty{1.7}{\percent} 范围之内。
\end{solution}

\begin{solution}{exo:b2:ore-deposits:5}
(a) 铂，$C_S = 10 \times 10^4 (R + 1)/(R + 10^4)$ ppb：$R = 100$ 时为 \qty{1.00}{ppm}，
\num{1000} 时为 \qty{9.10}{ppm}，\num{10000} 时为 \qty{50.0}{ppm}，\num{100000} 时为 \qty{90.9}{ppm}。
镍，$C_S = 200 \times 500 (R + 1)/(R + 500)$ ppm：\qty{1.68}{\percent}，\qty{6.67}{\percent}，
\qty{9.52}{\percent}，\qty{9.95}{\percent}。(b) $D\,C_L$：铂为 \qty{100}{ppm}，
镍为 \qty{10}{\percent}。(c) 自 $R = \num{10000}$ 起，镍的含量已在其上限的 \qty{5}{\percent} 以内，
$R$ 变化十倍时其含量几乎不变；铂在此处仍近似翻倍，
故其含量可以反映 $R$。
\end{solution}

\begin{solution}{exo:b2:ore-deposits:6}
(a) $\mu_M = 2.373 - 0.364 - 2 = 0.009$；$\sigma_M = 0.702$；中位数 $10^{0.009} = \qty{1.02}{Mt}$。
(b) $z = (1 - 0.009)/0.702 = 1.412$，$1 - \Phi(1.412) = 0.0790$：422 个\omterm{def:b2:ore-deposits:deposit}{矿床}中约 33 个。
(c) $\sigma = \sqrt{0.654^2 + 0.199^2} = 0.684$，$z = 1.449$，$1 - \Phi = 0.0737$，约 31 个：
弱相关性使大\omterm{def:b2:ore-deposits:deposit}{矿床}多出几个。
\end{solution}

\begin{solution}{exo:b2:ore-deposits:7}
全部块体的平均品位为 $0.40\,e^{0.18} = \qty{0.479}{\percent}$。对 $c = 0.3$：$z = \ln(0.40/0.3)/0.60 =
0.4795$，$\Phi(z) = 0.6842$，$\Phi(z + s) = 0.8598$：$T = \qty{684}{Mt}$，$\bar g = 0.479 \times
0.8598/0.6842 = \qty{0.602}{\percent}$，$M = 1000 \times 0.00479 \times 0.8598 = \qty{4.12}{Mt}$。
对 $c = 0.6$：$z = -0.6758$，$\Phi = 0.2496$，$\Phi(z + s) = 0.4698$：\qty{250}{Mt}，
品位 \qty{0.901}{\percent}，金属量 \qty{2.25}{Mt}。由图可见，\omterm{def:b2:ore-deposits:grade}{边界品位}为 \qty{0.3}{\percent} 时，蓝色曲线约为
\qty{680}{Mt}，红色曲线约为 \qty{0.6}{\percent}。
\end{solution}

\begin{solution}{exo:b2:ore-deposits:8}
(a) 流体：$\num{5e6}/(\num{1e-3} \times 0.5) = \qty{1.00e10}{t}$；\omterm{def:b2:magma-emplacement:magma}{岩浆}：$\num{1e10}/0.035 =
\qty{2.86e11}{t}$。(b) $\num{2.86e14}/2600 = \qty{1.10e11}{m^3}$，即 \qty{110}{km^3}，相当于
边长约 \qty{4.8}{km} 的\omterm{def:b2:magma-emplacement:magma}{岩浆}立方体。(c) \omterm{def:b2:magma-emplacement:magma}{岩浆}中含铜 $\num{2.86e11}
\times \num{50e-6} = \qty{1.43e7}{t}$ 铜，矿石所含铜占其中的 $5/14.3 = 0.350$。
\end{solution}

\begin{solution}{exo:b2:ore-deposits:9}
(a) 对数的均值为 1.411，标准差为 0.0664；中位数 $10^{1.411} = \qty{25.8}{ppm}$；
阈值 $10^{1.411 + 0.133} = 10^{1.544} = \qty{35.0}{ppm}$。(b) $(4 \times 400 + 36 \times 25.8)/40 =
\qty{63.2}{ppm}$，高于 \qty{35.0}{ppm}：经九份清洁
沉积物稀释后仍为异常。(c) $4 \times 400 + (A - 4) \times 25.8 = 35.0\,A$，由此得 $A = 4 \times (400 -
25.8)/(35.0 - 25.8) = \qty{163}{km^2}$。
\end{solution}

\begin{solution}{exo:b2:ore-deposits:10}
(a) $e^{su}\,e^{-u^2/2} = e^{-(u - s)^2/2 + s^2/2}$；除以 $\sqrt{2\pi}$。(b) $E[e^{\mu + su}] =
e^{\mu}\int e^{su}\varphi(u)\,\mathrm{d}u = e^{\mu + s^2/2}\int\varphi(u - s)\,\mathrm{d}u = e^{\mu +
s^2/2}$。(c) $\Phi$ 递增且 $s > 0$，故 $\Phi(z + s) > \Phi(z)$：保留的块体比舍弃的块体品位更高，
因此每个\omterm{def:b2:ore-deposits:grade}{边界品位}保留的金属份额都大于岩石份额；
二者之差随 $s$（\omterm{def:b2:ore-deposits:deposit}{矿床}的变异性）增大而增大。
\end{solution}

\begin{solution}{exo:b2:ore-deposits:11}
(a) 由 $D(R + 1)/(R + D) = fD$ 得 $R + 1 = fR + fD$，$R = (fD - 1)/(1 - f)$。(b) $f = 0.5$：$R = D - 2$，
得 298 与 \num{9998}；$f = 0.9$：$R = 9D - 10$，得 \num{2690} 与 \num{89990}。(c) 镍在
镍铜矿石的范围（100 至 \num{2000}）内即达到其上限的 \qty{50}{\percent}，略高于此时为 90；
铂需要 $R$ 为 \num{10000} 至 \num{100000}，即铂族元素礁的范围。
\end{solution}

\begin{solution}{exo:b2:ore-deposits:12}
(a) 由 $\ln M = m + s u$，分位数以上的\omterm{def:b2:ore-deposits:deposit}{矿床}满足 $u > \Phi^{-1}(q)$；其金属量为
$\int_{\Phi^{-1}(q)}^{\infty} e^{m + su}\varphi(u)\,\mathrm{d}u = e^{m + s^2/2}[1 - \Phi(\Phi^{-1}(q) -
s)]$，除以总量 $e^{m + s^2/2}$ 得 $\Phi(s - \Phi^{-1}(q))$。(b) $s = 0.702 \ln 10 =
1.616$；$\Phi(1.616 - 1.282) = \Phi(0.334) = 0.631$。(c) $\Phi(1.616 - 2.326) = \Phi(-0.710) = 0.239$，
与最大四个\omterm{def:b2:ore-deposits:deposit}{矿床}的 0.266 相比，最大\omterm{def:b2:ore-deposits:deposit}{矿床}比对数正态律
所预期的略大。
\end{solution}

\begin{solution}{pb:b2:ore-deposits:1}
\textbf{1.} \qty{330}{m}；$\sum l\,g = 138.4$，平均 $138.4/330 = \qty{0.419}{\percent}$。
\textbf{2.} $3.24/8 = \qty{0.405}{\percent}$；每个区间代表一段岩芯，
进而代表一段岩石体积：长区间须比短区间计入更多。
\textbf{3.} $230/330 = 0.697$，品位为 $118.6/230 = \qty{0.516}{\percent}$。
\textbf{4.} $4190/28 = 150$。
\textbf{5.} $8100/28 = 289$。
\textbf{6.} $0.38\,e^{0.55^2/2} = 0.38\,e^{0.1513} = \qty{0.442}{\percent}$；钻孔
品位为 \qty{0.419}{\percent}，与之相差在 \qty{6}{\percent} 以内：具有代表性。
\textbf{7.} $\pi \times 450^2 \times 700 = \qty{4.45e8}{m^3}$；$\times 2.7 = \num{1.20e9}$ 吨，即
\qty{1200}{Mt}。
\textbf{8.} $z = \ln(0.38/0.3)/0.55 = 0.4298$；$\Phi(z) = 0.6663$。
\textbf{9.} $z = 1.167$，0.4298，$-0.4990$；$\Phi = 0.8784$，0.6663，0.3089；$T = \qty{1050}{Mt}$，
\qty{800}{Mt}，\qty{371}{Mt}。
\textbf{10.} $\Phi(z + s) = \Phi(0.9798) = 0.8364$；$\bar g = 0.442 \times 0.8364/0.6663 =
\qty{0.555}{\percent}$。
\textbf{11.} $0.442 \times 0.9570/0.8784 = \qty{0.482}{\percent}$；$0.442 \times 0.5203/0.3089 =
\qty{0.744}{\percent}$。
\textbf{12.} 吨位 $\times\, 370.7/799.6 = 0.464$；品位 $\times\, 0.744/0.555 = 1.34$。
\textbf{13.} 在中位数处，\qty{0.38}{\percent}（$z = 0$）：\qty{600}{Mt}，品位 $0.442 \times
\Phi(0.55)/0.5 = 0.442 \times 0.7088/0.5 = \qty{0.627}{\percent}$。
\textbf{14.} $1200 \times 0.00442 = \qty{5.30}{Mt}$。
\textbf{15.} $(\log_{10} 1200 - 2.373)/0.654 = (3.0792 - 2.373)/0.654 = 1.080$；$\Phi = 0.8599$：
第 86 百分位，属大\omterm{def:b2:ore-deposits:deposit}{矿床}。
\textbf{16.} $(\log_{10} 0.442 + 0.364)/0.199 = (-0.3546 + 0.364)/0.199 = 0.047$；$\Phi = 0.5187$：
普通品位，第 52 百分位。
\textbf{17.} 考虑相关性时为 0.702，不考虑时为 0.684。
\textbf{18.} $\mu_M = 0.009$；$(\log_{10} 5.30 - 0.009)/0.702 = (0.7243 - 0.009)/0.702 = 1.019$；$\Phi =
0.8459$：第 85 百分位。
\textbf{19.} $422 \times (1 - 0.8459) = 65$。
\textbf{20.} $M = T_0\,e^{\mu + s^2/2}\,\Phi(z + s) = 1200 \times 0.00442 \times 0.9570 =
\qty{5.08}{Mt}$。
\textbf{21.} $1200 \times 0.00442 \times 0.5203 = \qty{2.76}{Mt}$。
\textbf{22.} $\Phi(z + s) = 0.8364$ 的铜，对应 $\Phi(z) = 0.6663$ 的吨位。
\textbf{23.} \omterm{def:b2:ore-deposits:grade}{边界品位}从 0.3 降至 \qty{0.2}{\percent}，吨位增加 $1054 - 800 = \qty{254}{Mt}$
（\qty{32}{\percent}），铜量增加 $5.08 - 4.44 = \qty{0.64}{Mt}$（\qty{14}{\percent}）：
新增块体的平均品位为 $0.64/254 = \qty{0.25}{\percent}$，为该\omterm{def:b2:ore-deposits:deposit}{矿床}中最贫的部分。
\textbf{24.} $1200 \times 0.00442 \times 0.8364 = 800 \times 0.00555$：该勘探区在 \qty{0.3}{\percent} 的\omterm{def:b2:ore-deposits:grade}{边界品位}以上含有
\textbf{\qty[text-series-to-math=true]{4.44}{Mt}} 铜。
\end{solution}
